\documentclass[11pt]{article}

\usepackage[margin=1in]{geometry}
\usepackage{amsmath,amssymb,bm}

\title{Linearized Barotropic Potential-Vorticity Equation in Mercator Coordinates}
\author{}
\date{}

\begin{document}

\maketitle

\section{Mercator Coordinates}

Let
\[
x=\lambda,\qquad 
y=\log \tan\left(\frac{\pi}{4}+\frac{\phi}{2}\right),
\qquad 
\frac{d\phi}{dy}=\cos\phi .
\]
In these Mercator coordinates, the spherical metric is
\[
ds^2=a^2\cos^2\phi\,(dx^2+dy^2),
\]
where \(a\) is the radius of the sphere.

The computational domain is taken to be all longitudes and latitudes from
\(80^\circ\mathrm{S}\) to \(80^\circ\mathrm{N}\):
\[
0\leq \lambda < 2\pi,
\qquad
-\frac{80\pi}{180}\leq \phi \leq \frac{80\pi}{180}.
\]
Equivalently,
\[
0\leq x<2\pi,
\qquad
y_S\leq y\leq y_N,
\]
where
\[
y_S
=
\log\tan\left(
\frac{\pi}{4}
-\frac{80\pi}{360}
\right),
\qquad
y_N
=
\log\tan\left(
\frac{\pi}{4}
+\frac{80\pi}{360}
\right).
\]
Since the Mercator coordinate is antisymmetric about the equator,
\[
y_S=-y_N .
\]

\section{Perturbation Stream Function and Vorticity}

Let \(\psi(x,y,t)\) be the perturbation stream function. The perturbation velocity components are
\[
u'=-\frac{1}{a\cos\phi}\psi_y,
\qquad
v'=\frac{1}{a\cos\phi}\psi_x .
\]

The perturbation relative vorticity is
\[
\zeta'
=
\nabla^2\psi
=
\frac{1}{a^2\cos^2\phi}
\left(\psi_{xx}+\psi_{yy}\right).
\]

Define the Mercator Laplacian
\[
\Delta_M=\partial_x^2+\partial_y^2 .
\]
Then
\[
\zeta'
=
\frac{1}{a^2\cos^2\phi}\Delta_M\psi .
\]

\section{Background Flow and Background PV}

The background flow is prescribed by the physical zonal and meridional winds
\[
U=U(\lambda,\phi),
\qquad
V=V(\lambda,\phi).
\]
In Mercator coordinates these are regarded as
\[
U=U(x,\phi(y)),
\qquad
V=V(x,\phi(y)).
\]

The Coriolis parameter is
\[
f=2\Omega\sin\phi .
\]

The background absolute vorticity is
\[
Q_0=f+\zeta_0,
\]
where the background relative vorticity is
\[
\zeta_0
=
\frac{1}{a\cos\phi}
\left[
V_x-\frac{1}{\cos\phi}\partial_y(U\cos\phi)
\right].
\]
Equivalently,
\[
\zeta_0
=
\frac{1}{a\cos\phi}
\left(
V_x-U_y+U\sin\phi
\right).
\]

The background material derivative is
\[
\mathcal D_0
=
\partial_t
+
\frac{U}{a\cos\phi}\partial_x
+
\frac{V}{a\cos\phi}\partial_y .
\]

Define the Jacobian
\[
J(A,B)=A_xB_y-A_yB_x .
\]

\section{Linearized Barotropic PV Equation}

The inviscid, unforced linearized barotropic potential-vorticity equation for the perturbation stream function is
\[
\mathcal D_0
\left[
\frac{1}{a^2\cos^2\phi}
\Delta_M\psi
\right]
+
\frac{1}{a^2\cos^2\phi}
J(\psi,Q_0)
=0 .
\]

Now include linear friction proportional to minus the perturbation relative vorticity,
\[
-r\nabla^2\psi=-r\zeta',
\]
and include an external forcing \(F(x,y,t)\). The forced-damped linearized equation is
\[
\boxed{
\mathcal D_0
\left[
\frac{1}{a^2\cos^2\phi}
\Delta_M\psi
\right]
+
\frac{1}{a^2\cos^2\phi}
J(\psi,Q_0)
=
-r
\left[
\frac{1}{a^2\cos^2\phi}
\Delta_M\psi
\right]
+
F .
}
\]

Equivalently,
\[
\boxed{
\left(
\partial_t
+
\frac{U}{a\cos\phi}\partial_x
+
\frac{V}{a\cos\phi}\partial_y
\right)
\left[
\frac{1}{a^2\cos^2\phi}
\left(\psi_{xx}+\psi_{yy}\right)
\right]
+
\frac{1}{a^2\cos^2\phi}
\left(
\psi_x Q_{0y}-\psi_y Q_{0x}
\right)
=
-r
\left[
\frac{1}{a^2\cos^2\phi}
\left(\psi_{xx}+\psi_{yy}\right)
\right]
+
F .
}
\]

If the background flow is not itself a steady solution of the nonlinear barotropic
potential-vorticity equation, then an additional inhomogeneous term remains:
\[
\boxed{
\mathcal D_0
\left[
\frac{1}{a^2\cos^2\phi}
\Delta_M\psi
\right]
+
\frac{1}{a^2\cos^2\phi}
J(\psi,Q_0)
=
-r
\left[
\frac{1}{a^2\cos^2\phi}
\Delta_M\psi
\right]
+
F
-
\mathcal D_0 Q_0 .
}
\]

\section{Gaussian Forcing in the Ni\~no 3.4 Region}

The external forcing \(F\) is chosen to be a Gaussian function of longitude
and latitude, localized in the Ni\~no 3.4 region. The Ni\~no 3.4 box is
\[
5^\circ\mathrm{S}\leq \phi \leq 5^\circ\mathrm{N},
\qquad
170^\circ\mathrm{W}\leq \lambda \leq 120^\circ\mathrm{W}.
\]
Since the longitude domain is written as \(0\leq \lambda<2\pi\), the same longitude
range is
\[
190^\circ\mathrm{E}\leq \lambda \leq 240^\circ\mathrm{E}.
\]
In radians,
\[
-\frac{5\pi}{180}\leq \phi \leq \frac{5\pi}{180},
\qquad
\frac{190\pi}{180}\leq \lambda \leq \frac{240\pi}{180}.
\]

The center of the box is
\[
\phi_c=0,
\qquad
\lambda_c=\frac{215\pi}{180}.
\]
Equivalently, \(\lambda_c=145^\circ\mathrm{W}=215^\circ\mathrm{E}\).

Define the Gaussian widths
\[
\sigma_\lambda=\frac{25\pi}{180},
\qquad
\sigma_\phi=\frac{5\pi}{180}.
\]

Because longitude is periodic, define the wrapped longitudinal distance
\[
d_\lambda(\lambda,\lambda_c)
=
\operatorname{atan2}
\left(
\sin(\lambda-\lambda_c),
\cos(\lambda-\lambda_c)
\right).
\]
This gives the shortest signed angular distance between \(\lambda\) and
\(\lambda_c\).

The forcing is prescribed as
\[
\boxed{
F(\lambda,\phi,t)
=
F_0\,S(t)
\exp
\left[
-\frac{1}{2}
\left(
\frac{d_\lambda(\lambda,\lambda_c)^2}{\sigma_\lambda^2}
+
\frac{(\phi-\phi_c)^2}{\sigma_\phi^2}
\right)
\right].
}
\]

Here \(F_0\) is the forcing amplitude and \(S(t)\) is a nondimensional temporal
envelope. For a steady forcing,
\[
S(t)=1.
\]
For an oscillatory forcing, one may take
\[
S(t)=\cos(\omega t).
\]
For a pulsed forcing, one may take
\[
S(t)
=
\exp\left[
-\frac{(t-t_c)^2}{2\sigma_t^2}
\right].
\]

In Mercator coordinates, since \(x=\lambda\) and \(\phi=\phi(y)\), the forcing is
\[
\boxed{
F(x,y,t)
=
F_0\,S(t)
\exp
\left[
-\frac{1}{2}
\left(
\frac{d_x(x,x_c)^2}{\sigma_\lambda^2}
+
\frac{(\phi(y)-\phi_c)^2}{\sigma_\phi^2}
\right)
\right],
}
\]
where
\[
x_c=\lambda_c=\frac{215\pi}{180},
\]
and
\[
d_x(x,x_c)
=
\operatorname{atan2}
\left(
\sin(x-x_c),
\cos(x-x_c)
\right).
\]

\section{Prognostic-Diagnostic Form}

Define the perturbation potential vorticity, equivalently the perturbation
absolute-vorticity anomaly in the barotropic constant-depth case, by
\[
q
=
\zeta'
=
\frac{1}{a^2\cos^2\phi}\Delta_M\psi .
\]

Then the forced-damped linearized equation becomes the prognostic equation
\[
\boxed{
\mathcal D_0 q
+
\frac{1}{a^2\cos^2\phi}J(\psi,Q_0)
=
-rq+F .
}
\]

Equivalently,
\[
\boxed{
q_t
+
\frac{U}{a\cos\phi}q_x
+
\frac{V}{a\cos\phi}q_y
+
\frac{1}{a^2\cos^2\phi}
\left(
\psi_x Q_{0y}-\psi_y Q_{0x}
\right)
=
-rq+F .
}
\]

With the Ni\~no 3.4 Gaussian forcing, this becomes
\[
\boxed{
q_t
+
\frac{U}{a\cos\phi}q_x
+
\frac{V}{a\cos\phi}q_y
+
\frac{1}{a^2\cos^2\phi}
\left(
\psi_x Q_{0y}-\psi_y Q_{0x}
\right)
=
-rq
+
F_0\,S(t)
\exp
\left[
-\frac{1}{2}
\left(
\frac{d_x(x,x_c)^2}{\sigma_\lambda^2}
+
\frac{(\phi(y)-\phi_c)^2}{\sigma_\phi^2}
\right)
\right].
}
\]

The diagnostic relation between \(q\) and \(\psi\) is the elliptic equation
\[
\boxed{
\frac{1}{a^2\cos^2\phi}\Delta_M\psi=q .
}
\]
Equivalently,
\[
\boxed{
\Delta_M\psi
=
a^2\cos^2\phi\,q .
}
\]

\section{Two-Step Time Advancement}

Suppose \(q^n\) and \(\psi^n\) are known at time \(t^n\). A simple explicit
time discretization from \(t^n\) to \(t^{n+1}=t^n+\Delta t\) is as follows.

\subsection{Step 1: Advance the PV}

Use the prognostic equation to advance \(q\):
\[
\boxed{
\frac{q^{n+1}-q^n}{\Delta t}
+
\frac{U}{a\cos\phi}q_x^n
+
\frac{V}{a\cos\phi}q_y^n
+
\frac{1}{a^2\cos^2\phi}
J(\psi^n,Q_0)
=
-rq^n+F^n .
}
\]

Thus,
\[
\boxed{
q^{n+1}
=
q^n
-\Delta t
\left[
\frac{U}{a\cos\phi}q_x^n
+
\frac{V}{a\cos\phi}q_y^n
+
\frac{1}{a^2\cos^2\phi}
J(\psi^n,Q_0)
\right]
+
\Delta t
\left[
-rq^n+F^n
\right].
}
\]

For the Ni\~no 3.4 Gaussian forcing,
\[
\boxed{
F^n(x,y)
=
F_0\,S(t^n)
\exp
\left[
-\frac{1}{2}
\left(
\frac{d_x(x,x_c)^2}{\sigma_\lambda^2}
+
\frac{(\phi(y)-\phi_c)^2}{\sigma_\phi^2}
\right)
\right].
}
\]

If the background is not a steady solution, replace the right-hand side by
\[
-rq^n+F^n-\mathcal D_0 Q_0 .
\]

\subsection{Step 2: Solve the Diagnostic Equation for the Stream Function}

After \(q^{n+1}\) has been computed, recover \(\psi^{n+1}\) by solving
\[
\boxed{
\frac{1}{a^2\cos^2\phi}\Delta_M\psi^{n+1}=q^{n+1}.
}
\]
Equivalently,
\[
\boxed{
\Delta_M\psi^{n+1}
=
a^2\cos^2\phi\,q^{n+1}.
}
\]

Thus, at each time step, one first updates the prognostic PV variable \(q\), and
then solves the elliptic inversion problem for the stream function \(\psi\).

\section{Finite-Difference Formulation in Space and Time}

Let the computational grid be
\[
x_i=i\Delta x,
\qquad
i=0,\ldots,N_x-1,
\qquad
\Delta x=\frac{2\pi}{N_x},
\]
and
\[
y_j=y_S+j\Delta y,
\qquad
j=0,\ldots,N_y-1,
\qquad
\Delta y=\frac{y_N-y_S}{N_y-1}.
\]
The time grid is
\[
t^n=t^0+n\Delta t .
\]

Write
\[
q_{i,j}^n \approx q(x_i,y_j,t^n),
\qquad
\psi_{i,j}^n \approx \psi(x_i,y_j,t^n).
\]
Define
\[
\phi_j=\phi(y_j),
\qquad
m_j=a\cos\phi_j .
\]
Thus
\[
\frac{1}{a\cos\phi_j}=\frac{1}{m_j},
\qquad
\frac{1}{a^2\cos^2\phi_j}=\frac{1}{m_j^2}.
\]

The longitude direction is periodic. Therefore, for any grid field \(A_{i,j}\),
\[
A_{i+N_x,j}=A_{i,j}.
\]
In particular,
\[
A_{-1,j}=A_{N_x-1,j},
\qquad
A_{N_x,j}=A_{0,j}.
\]

The centered finite-difference approximations are
\[
\delta_x A_{i,j}
=
\frac{A_{i+1,j}-A_{i-1,j}}{2\Delta x},
\qquad
\delta_y A_{i,j}
=
\frac{A_{i,j+1}-A_{i,j-1}}{2\Delta y}.
\]
The second derivatives are
\[
\delta_{xx}A_{i,j}
=
\frac{A_{i+1,j}-2A_{i,j}+A_{i-1,j}}{\Delta x^2},
\qquad
\delta_{yy}A_{i,j}
=
\frac{A_{i,j+1}-2A_{i,j}+A_{i,j-1}}{\Delta y^2}.
\]

The discrete Mercator Laplacian is
\[
\Delta_M^h A_{i,j}
=
\delta_{xx}A_{i,j}
+
\delta_{yy}A_{i,j}.
\]

A centered approximation to the Jacobian is
\[
J^h(A,B)_{i,j}
=
(\delta_x A)_{i,j}(\delta_y B)_{i,j}
-
(\delta_y A)_{i,j}(\delta_x B)_{i,j}.
\]

The discrete diagnostic relation between \(q\) and \(\psi\) is
\[
\boxed{
q_{i,j}^n
=
\frac{1}{m_j^2}
\Delta_M^h \psi_{i,j}^n .
}
\]
Equivalently,
\[
\boxed{
\Delta_M^h \psi_{i,j}^n
=
m_j^2 q_{i,j}^n .
}
\]

Written out explicitly,
\[
\boxed{
\frac{\psi_{i+1,j}^n-2\psi_{i,j}^n+\psi_{i-1,j}^n}{\Delta x^2}
+
\frac{\psi_{i,j+1}^n-2\psi_{i,j}^n+\psi_{i,j-1}^n}{\Delta y^2}
=
m_j^2 q_{i,j}^n .
}
\]

The background fields are evaluated on the grid as
\[
U_{i,j}=U(x_i,\phi_j),
\qquad
V_{i,j}=V(x_i,\phi_j),
\qquad
Q_{0,i,j}=Q_0(x_i,y_j).
\]

The Gaussian forcing is evaluated on the grid as
\[
\boxed{
F_{i,j}^n
=
F_0\,S(t^n)
\exp
\left[
-\frac{1}{2}
\left(
\frac{d_x(x_i,x_c)^2}{\sigma_\lambda^2}
+
\frac{(\phi_j-\phi_c)^2}{\sigma_\phi^2}
\right)
\right].
}
\]
Here
\[
\phi_c=0,
\qquad
x_c=\frac{215\pi}{180},
\qquad
\sigma_\lambda=\frac{25\pi}{180},
\qquad
\sigma_\phi=\frac{5\pi}{180},
\]
and
\[
d_x(x_i,x_c)
=
\operatorname{atan2}
\left(
\sin(x_i-x_c),
\cos(x_i-x_c)
\right).
\]

Using forward Euler in time and centered differences in space, the fully discrete
PV update is
\[
\boxed{
\begin{aligned}
\frac{q_{i,j}^{n+1}-q_{i,j}^{n}}{\Delta t}
&+
\frac{U_{i,j}}{m_j}
(\delta_x q^n)_{i,j}
+
\frac{V_{i,j}}{m_j}
(\delta_y q^n)_{i,j}
+
\frac{1}{m_j^2}
J^h(\psi^n,Q_0)_{i,j}
\\
&=
-r q_{i,j}^{n}
+
F_{i,j}^{n}.
\end{aligned}
}
\]

Therefore,
\[
\boxed{
\begin{aligned}
q_{i,j}^{n+1}
=
q_{i,j}^{n}
&-
\Delta t
\left[
\frac{U_{i,j}}{m_j}
(\delta_x q^n)_{i,j}
+
\frac{V_{i,j}}{m_j}
(\delta_y q^n)_{i,j}
+
\frac{1}{m_j^2}
J^h(\psi^n,Q_0)_{i,j}
\right]
\\
&+
\Delta t
\left[
-r q_{i,j}^{n}
+
F_{i,j}^{n}
\right].
\end{aligned}
}
\]

Written out without operator notation, this is
\[
\boxed{
\begin{aligned}
q_{i,j}^{n+1}
=
q_{i,j}^{n}
&-
\Delta t
\left[
\frac{U_{i,j}}{m_j}
\frac{q_{i+1,j}^{n}-q_{i-1,j}^{n}}{2\Delta x}
+
\frac{V_{i,j}}{m_j}
\frac{q_{i,j+1}^{n}-q_{i,j-1}^{n}}{2\Delta y}
\right]
\\
&-
\Delta t
\left[
\frac{1}{m_j^2}
\left(
\frac{\psi_{i+1,j}^{n}-\psi_{i-1,j}^{n}}{2\Delta x}
\frac{Q_{0,i,j+1}-Q_{0,i,j-1}}{2\Delta y}
\right.\right.
\\
&\hspace{4.5cm}
\left.\left.
-
\frac{\psi_{i,j+1}^{n}-\psi_{i,j-1}^{n}}{2\Delta y}
\frac{Q_{0,i+1,j}-Q_{0,i-1,j}}{2\Delta x}
\right)
\right]
\\
&+
\Delta t
\left[
-r q_{i,j}^{n}
+
F_0\,S(t^n)
\exp
\left[
-\frac{1}{2}
\left(
\frac{d_x(x_i,x_c)^2}{\sigma_\lambda^2}
+
\frac{(\phi_j-\phi_c)^2}{\sigma_\phi^2}
\right)
\right]
\right].
\end{aligned}
}
\]

After this prognostic update, the stream function at the new time level is obtained
from the discrete elliptic inversion
\[
\boxed{
\Delta_M^h \psi_{i,j}^{n+1}
=
m_j^2 q_{i,j}^{n+1}.
}
\]
That is,
\[
\boxed{
\frac{\psi_{i+1,j}^{n+1}-2\psi_{i,j}^{n+1}+\psi_{i-1,j}^{n+1}}{\Delta x^2}
+
\frac{\psi_{i,j+1}^{n+1}-2\psi_{i,j}^{n+1}+\psi_{i,j-1}^{n+1}}{\Delta y^2}
=
m_j^2 q_{i,j}^{n+1}.
}
\]

This is a linear elliptic problem for the unknowns \(\psi_{i,j}^{n+1}\).

\section{Boundary Conditions}

The longitude direction is periodic:
\[
q(0,y,t)=q(2\pi,y,t),
\qquad
\psi(0,y,t)=\psi(2\pi,y,t).
\]
Equivalently, on the discrete grid,
\[
q_{i+N_x,j}^n=q_{i,j}^n,
\qquad
\psi_{i+N_x,j}^n=\psi_{i,j}^n.
\]

At the artificial northern and southern boundaries,
\[
\phi=\pm 80^\circ,
\]
one must impose a meridional boundary condition. Common choices are homogeneous
Dirichlet conditions,
\[
\psi=0
\qquad
\text{at}
\qquad
\phi=\pm 80^\circ,
\]
or homogeneous Neumann conditions,
\[
\frac{\partial \psi}{\partial y}=0
\qquad
\text{at}
\qquad
\phi=\pm 80^\circ.
\]

For homogeneous Dirichlet conditions, the discrete boundary conditions are
\[
\psi_{i,0}^{n+1}=0,
\qquad
\psi_{i,N_y-1}^{n+1}=0.
\]

For homogeneous Neumann conditions, one may impose
\[
\frac{\psi_{i,1}^{n+1}-\psi_{i,0}^{n+1}}{\Delta y}=0,
\qquad
\frac{\psi_{i,N_y-1}^{n+1}-\psi_{i,N_y-2}^{n+1}}{\Delta y}=0.
\]

The corresponding treatment of \(q\) at the meridional boundaries should be chosen
consistently with the elliptic inversion for \(\psi\).


\end{document}
